\documentclass[10pt,a4paper]{amsart}
\usepackage[margin=19mm]{geometry}
\usepackage{amsmath,amssymb,mathtools}
\usepackage[scr=rsfs,cal=euler]{mathalfa}
\usepackage{xfrac}
\usepackage[unicode,hidelinks,bookmarks=false]{hyperref}
\newcommand{\ie}{i.e.\ }
\newcommand{\cf}{cf.\ }
\newcommand{\resp}{respectively\ }
\newcommand{\Subsection}{\subsection}
\providecommand{\nobreakdash}{\nobreak\hbox{-}}
\setlength{\emergencystretch}{2em}
\multlinegap0pt
\usepackage{bm}
\usepackage{bbm}
\usepackage{dsfont}
\usepackage{xspace}
\usepackage{cancel}
\usepackage{mathtools}
\usepackage[shortlabels]{enumitem}
\usepackage{amsthm}
\makeatletter
\@ifundefined{thm}{\newtheorem{thm}{Thm}}{}
\@ifundefined{cor}{\newtheorem{cor}[thm]{Corollary}}{}
\@ifundefined{lem}{\newtheorem{lem}[thm]{Lemma}}{}
\makeatother
\DeclareMathOperator{\lcm}{lcm}
\title[The Combinatorial Nullstellensatz, Chevalley-Warning Theorem]{The Combinatorial Nullstellensatz, Chevalley-Warning Theorem and weak Finitesatz in skew polynomial rings}
\author{Gil Alon, Angelot Behajaina, Elad Paran}
\date{Source: arXiv:2508.07257v2 (CC BY 4.0)}
\begin{document}
\maketitle

\noindent\emph{Source and licence.} The Combinatorial Nullstellensatz, Chevalley-Warning Theorem and weak Finitesatz in skew polynomial rings. Version arXiv:2508.07257v2 (CC BY 4.0), distributed under the Creative Commons Attribution 4.0 International licence, as recorded in the arXiv licence metadata for this exact version and shown on the abstract page https://arxiv.org/abs/2508.07257v2. Full rights notice: licenses/arxiv-2508.07257-CC-BY-4.0.txt.\par\smallskip

\noindent\textit{Editorial scope.} Counted unit: the Lem whose source label is \texttt{lem:nullvanismalldeg} in the article (source0.tex lines 602--610, source label \texttt{lem:nullvanismalldeg}), delivered together with the article's own complete original proof of it (source0.tex lines 611--651, one \texttt{proof} environment of 41 lines). No mathematics is written, repaired or re-derived: statement and proof are the article's own text line for line. Required context retained uncounted: cor:evalformon (lines 302--304); lem:lcmfinifield (lines 593--598). Every definition, display and label of those lines is retained unchanged. Omitted from this package and not redistributed: 1--301, 305--592, 599--601, 652--709. Nothing inside a retained excerpt is omitted or altered: every retained body is delivered exactly as the article's lines are written, so no omission receipt of any kind accompanies this package. Citations: the retained excerpts carry no citation command at all, so no bibliography entry of the article is transcribed and no reference is redistributed. Reference adaptation: none was needed and none was made; every reference inside the delivered unit resolves inside it, so no unresolved reference or citation is produced. Printed slips kept literal: no printed word, symbol, hypothesis or number of the counted statement or of its proof was altered. Layout and class: the article's own class is \texttt{amsart} with packages including amsthm, amsmath, amssymb, graphicx, float, hyperref, mathtools, relsize, tikz, enumerate, appendix, pdflscape, multirow, adjustbox; the wrapper is the lane's pinned \texttt{amsart} editorial wrapper at 10pt on A4, which restates the theorem-like environments the retained text needs and adds the article's own macro definitions that the retained text needs (\texttt{\textbackslash lcm}), with extra wrapper packages (bm, bbm, dsfont, xspace, cancel, mathtools, enumitem). The delivered numbering is the unit's own. Assets: no retained span contains a figure environment, a graphics-inclusion command, a PDF-page inclusion, a table or a TikZ picture; no image file is copied into this package, the package declares no asset dependency and no image file is redistributed with it. Licence: version arXiv:2508.07257v2 of this article is distributed under the Creative Commons Attribution 4.0 International licence, as recorded in the arXiv licence metadata for this exact version and shown on the abstract page https://arxiv.org/abs/2508.07257v2. A full rights notice is delivered at licenses/arxiv-2508.07257-CC-BY-4.0.txt. Characters: every retained excerpt is pure ASCII, so the delivered engine, XeLaTeX with its default Latin Modern text font, reports no missing character.
\begin{cor}\label{cor:evalformon}
Let $f \in R$ and ${\bf a} \in D^{n,\sigma}$. For any ordering of the variables, the corresponding ordered evaluation is the unique $\ell \in D$ such that $f-\ell \in \mathfrak{m}_{{\bf a}}$. In particular, the ordered evaluation is independent of the order.
\end{cor}

\begin{lem}\label{lem:lcmfinifield}
We have
$$
\lcm(x-a \mid a \in F)=x^{\frac{m}{\theta}(p^\theta-1)+1}-x.
$$
\end{lem}

\begin{lem}\label{lem:nullvanismalldeg}
Let $h \in \mathscr{I}(F^{n,\sigma})$ be such that:
\begin{enumerate}[i)]
\item
$\deg_{x_{i}}(h) \leq \frac{m}{\theta}(p^\theta-1)$, for all $1 \leq i \leq n$; \label{eq1:lem:nullvanismalldeg}
\item for every $  1 \leq j<i\leq n$, no monomial in $h$ is divisible by $x_ix_j^{p^\theta}$.\label{eq2:lem:nullvanismalldeg}
\end{enumerate}
Then $h=0$.
\end{lem}

\begin{proof}
We prove the result by induction on $n \geq 1$. 

For $n=1$, let $h \in \mathscr{I}(F)$ satisfy \eqref{eq1:lem:nullvanismalldeg} and \eqref{eq2:lem:nullvanismalldeg}. By Lemma \ref{lem:lcmfinifield}, $x_1^{\frac{m}{\theta}(p^\theta-1)+1}-x_1$ divides $h$. Then, from \eqref{eq1:lem:nullvanismalldeg}, we obtain $h=0$. Therefore the theorem holds for $n=1$

Now, assume that the theorem holds for all integers $1,\dots,n-1$ with $n \geq 2$. Let $h \in \mathscr{I}(F^{n,\sigma})$ satisfy \eqref{eq1:lem:nullvanismalldeg} and \eqref{eq2:lem:nullvanismalldeg}. Assume to the contrary that $h \neq 0$. Write 
$$h=x_1^e h_{e}(x_2,\dots,x_n)+x_1^{e-1} h_{e-1}(x_2,\dots,x_n)+\cdots+x_1 h_{1}(x_2,\dots,x_n)+ h_{0}(x_2,\dots,x_n),
$$ where $h_i(x_2,\dots,x_n) \in F[x_2,\dots,x_n;\sigma]$ ($1 \leq i \leq e$) and $h_e(x_2,\dots,x_n)\neq 0$. We consider two cases.
\vskip 1mm
\noindent
$\bullet$ {\bf First suppose that $e\leq p^\theta-1$}. Let $\lambda \in \mathscr{W}_F$. Then, for every ${\bf v}=\omega_\lambda (a_1,\dots,a_n)\in \mathfrak{S}_{\lambda}^n= \omega_\lambda K^n$, using Corollary \ref{cor:evalformon},
\begin{equation}\label{eq:hvv}
\begin{aligned}
0=h({\bf v})
&={\rm N}^\sigma_e(\omega_\lambda a_1)
  \sigma^e\big(h_e(\omega_\lambda(a_2,\dots,a_n))\big)\\
&\quad+\cdots+{\rm N}^\sigma_1(\omega_\lambda a_1)
  \sigma\big(h_1(\omega_\lambda(a_2,\dots,a_n))\big)
  +h_0(\omega_\lambda(a_2,\dots,a_n))\\
&=a_1^e {\rm N}^\sigma_e(\omega_\lambda)
  \sigma^e\big(h_e(\omega_\lambda(a_2,\dots,a_n))\big)\\
&\quad+\cdots+a_1{\rm N}^\sigma_1(\omega_\lambda)
  \sigma\big(h_1(\omega_\lambda(a_2,\dots,a_n))\big)
  +h_0(\omega_\lambda(a_2,\dots,a_n)).
\end{aligned}
\end{equation} Fix $(a_2,\dots,a_n) \in K^{n-1}$. Since $e \leq p^\theta-1$, ${\rm N}^\sigma_e(\omega_\lambda) \neq 0$ and \eqref{eq:hvv} holds for all $a_1 \in K$, we have $h_e(\omega_{\lambda}(a_2,\dots,a_n))=0$. Hence, for every $\lambda \in \mathscr{W}_F$, the polynomial $h_e \in F[x_2,\dots,x_n;\sigma]$ vanishes on $\mathfrak{S}_{\lambda}^{n-1}$. Therefore $h_e$ also vanishes on $F^{n-1,\sigma}=\bigcup_{\lambda \in \mathscr{W}_F} \mathfrak{S}_{\lambda}^{n-1}$. Since $h_e$ also satisfies \eqref{eq1:lem:nullvanismalldeg} and \eqref{eq2:lem:nullvanismalldeg} for all $2 \leq i,j \leq n$, the induction hypothesis implies that $h_e=0$, a contradiction.
\vskip 1mm
\noindent
$\bullet$ {\bf Now suppose that $e>p^\theta-1$}. In this case, we have $c=h_e(x_2,\dots,x_n) \in F^*$; otherwise $h$ would contain $x_jx_1^{p^\theta}$ for some $j>1$, contradicting \eqref{eq2:lem:nullvanismalldeg}. Note that
\begin{align*}
0=h(\omega_\lambda(a_1,0,\dots,0))
&=\sigma^e(c){\rm N}^\sigma_e(\omega_\lambda a_1)\\
&\quad+{\rm N}^\sigma_{e-1}(\omega_\lambda a_1)
  \sigma^{e-1}\big(h_{e-1}({\bf 0})\big)+\cdots\\
&\quad+{\rm N}^\sigma_1(\omega_\lambda a_1)
  \sigma\big(h_1({\bf 0})\big)+h_0({\bf 0}),
\end{align*} for all $\lambda \in \mathscr{W}_{F}$ and all $a_1 \in K$. Letting 
$$g=\sigma^e(c)x_1^e+\sigma^{e-1}\big(h_{e-1}({\bf 0})\big)x_1^{e-1}+\cdots+\sigma\big(h_1({\bf 0})\big)x_1+h_0({\bf 0}) \in F[x_1;\sigma],$$ we see that $g$ vanishes on $F$. By Lemma \ref{lem:lcmfinifield}, $x_1^{\frac{m}{\theta}(p^\theta-1)+1}-x_1$ divides $g$, so $e \geq \frac{m}{\theta}(p^\theta-1)+1$, contradicting \eqref{eq1:lem:nullvanismalldeg}.

In both cases, we reach a contradiction. Consequently, we must have $h=0$, completing the proof by induction.
\end{proof}

\end{document}
